\chapter{Tests}
\label{ch:tests}
Both the agent have been tested locally during the building process continously,
to find bugs, test different configurations, and improve them. 

Extra attention was dedicated on checking the structure of the messages and actions proposed by the LLM,
since we cannot blindly rely on a non-deterministic way of extracting intentions from the Admin prompt.

To facilitate testing on multiple scenarios we decided to build a flexible map-builder tool 
(\texttt{run\_map\_builder.js}, served at \texttt{localhost:3000} through \texttt{npm start}).
After manual test, we implemented automated unit and coverage tests, which guarantees that the 
current behaviour remains consistent.
This is done via the built-in node test system, and code actions to verify that new code changes 
comply with the tests.
%The PDDL layer in particular is covered, verifying crate-ID propagation in \texttt{compileProblemPddl} 
%together with local obstacle-pushing and re-routing around crates that fail to be solved.


\subsection{Considerations}
\label{sub:res_and_lim}
We have found creating an appropriate, logically sound, system prompt for the
LLM not trivia, as it frequently tried to vier off the guardrails.

Some known limitations we found were: the occasional inability for the LLM to reason about associating 
multiple rewards to different tasks, resulting in missed opportunities and unhandled requests; not performing
the actions in the same sequential order as they were provided, leading to out of order answers.